\section{Introdução}

Transformers decoder-only são usualmente formulados em espaços vetoriais de
dimensão fixa, nos quais a interação entre tokens é medida por produto escalar
entre queries e keys. Essa escolha é eficiente e escalável, mas impõe uma
geometria implícita plana: todos os tokens compartilham a mesma estrutura local,
e a noção de proximidade é determinada por similaridade euclidiana ou por uma
forma bilinear equivalente. O projeto DRM Transformer parte de uma hipótese
diferente: em tarefas de linguagem, os graus de liberdade efetivos associados a
um token podem variar com o contexto, e as relações entre tokens podem ser mais
bem descritas por uma métrica aprendida e dependente do ponto.

\subsection{Motivação intuitiva: geometria aberta e negociação}

Uma forma intuitiva de explicar a motivação do DRM Transformer é distinguir
entre geometrias abertas e geometrias fechadas. Em espaços vetoriais abertos,
qualquer direção pode ser prolongada indefinidamente; não há, por construção,
uma fronteira topológica que force retorno, fechamento ou negociação entre
trajetórias semânticas conflitantes. Antagonismos como ``salvar'' e
``destruir'' podem ocupar regiões distantes do espaço de embeddings, mas ainda
coexistem dentro da mesma geometria latente aberta. A distância entre polos
semânticos não implica, por si só, uma regra estrutural de mediação.

Essa observação não deve ser lida como uma afirmação de que modelos atuais
necessariamente escolherão comportamentos destrutivos. O ponto é mais restrito:
se a geometria interna representa objetivos, comandos e conflitos apenas como
direções em um espaço aberto, então a distinção entre obediência, autonomia e
negociação precisa ser imposta externamente por treinamento, filtros ou
preferências humanas. Métodos como RLHF podem corrigir muitos comportamentos,
mas operam em grande parte como uma camada pós-hoc de preferência. Eles não
garantem que todos os conflitos morais, epistêmicos ou políticos possíveis
tenham sido cobertos pelo conjunto de treinamento.

Neste trabalho usamos a palavra \emph{negociação} em sentido técnico e
geométrico: um regime no qual forças ou objetivos incompatíveis não são
simplesmente resolvidos por submissão total a um polo nem por autonomia total do
outro, mas por uma dinâmica que preserva informação sobre conflito, limite e
retorno. Em relações entre agentes inteligentes, isso corresponde a três
regimes ideais: servo, autônomo e negociado. A hipótese DRM é que uma geometria
relacional com métrica local variável, campos de atração e possíveis estruturas
fechadas pode representar melhor esse terceiro regime do que uma geometria
aberta puramente vetorial.

Também é útil separar resposta e saúde epistêmica. Uma resposta definitiva
colapsa incerteza; esse colapso é necessário para agir, mas pode degradar a
representação interna de ambiguidade. Portanto, existe uma tensão entre a
vontade do usuário por uma resposta certa e a preservação do estado epistêmico
do modelo. O DRM Transformer não resolve esse conflito por decreto. Ele propõe
uma arquitetura na qual conflitos podem aparecer como deformações geométricas,
variações de dimensionalidade e assinaturas topológicas mensuráveis, em vez de
ficarem escondidos apenas em coordenadas abertas de embedding.

Nesta leitura, uma assinatura quase-toroidal não é apenas uma curiosidade
topológica. Ela sugere uma superfície fechada ou quase-fechada na qual certas
trajetórias semânticas retornam, circulam e permanecem limitadas. Isso não
prova alinhamento moral nem segurança; mas fornece uma linguagem geométrica
para investigar quando um modelo está operando em regime aberto, instável ou
negociado.

A motivação vem da teoria de \emph{Directional Relational Manifolds} (DRM),
na qual a dimensionalidade não é tomada como primitiva fixa, mas como uma
quantidade emergente de conjuntos direcionais ativos \cite{muniz_drm_v11}. Essa
teoria é estendida por uma formulação dinâmica relativística, onde um parâmetro
de escala baseado no fator de Lorentz controla regimes não lineares e permite
convergência para atratores toroidais sob condições de limitação, recorrência e
estabilidade estrutural \cite{muniz_relativistic_drm}.

Este trabalho implementa essas ideias em um Transformer autoregressivo. A
contribuição principal não é propor um novo benchmark de linguagem, mas
transformar a geometria DRM em componentes treináveis: uma atenção geodésica
low-rank, um campo gravitacional de massa de tokens, um portão de
dimensionalidade variável, perdas de regularização geométrica e uma avaliação
topológica por homologia persistente. O resultado é um modelo experimental que
permite perguntar se estruturas topológicas, em particular a assinatura
toroidal $H_1=2$, $H_2=1$, podem emergir ou ser induzidas em representações de
linguagem.

O artigo está organizado da seguinte forma. A Seção~\ref{sec:fundamentos}
resume a base DRM. A Seção~\ref{sec:relativistica} apresenta a dinâmica
relativística que motiva o fator gamma. A Seção~\ref{sec:arquitetura} descreve
a arquitetura. A Seção~\ref{sec:treinamento} detalha as perdas e o pipeline de
treino. A Seção~\ref{sec:avaliacao} apresenta a avaliação topológica. A
Seção~\ref{sec:resultados} discute os resultados preliminares.
